\documentclass[a4paper,11pt]{ltjsarticle}

\usepackage[margin=20mm,headheight=18pt]{geometry}
\usepackage{fontspec}
\usepackage{luatexja-fontspec}
\usepackage{xcolor}
\usepackage{listings}
\usepackage{booktabs}
\usepackage{tabularx}
\usepackage{enumitem}
\usepackage{fancyhdr}
\usepackage{titlesec}
\usepackage[hidelinks]{hyperref}
\usepackage{tikz}
\usetikzlibrary{arrows.meta,positioning}

\setmainfont{Noto Serif CJK JP}
\setsansfont{Noto Sans CJK JP}
\setmainjfont{Noto Serif CJK JP}
\setsansjfont{Noto Sans CJK JP}
\setmonofont{DejaVu Sans Mono}

\definecolor{Ink}{HTML}{17233C}
\definecolor{Blue}{HTML}{246BCE}
\definecolor{LightBlue}{HTML}{EAF2FD}
\definecolor{Orange}{HTML}{E87824}
\definecolor{SoftGray}{HTML}{F3F5F7}
\definecolor{CodeGray}{HTML}{4B5563}

\hypersetup{
  colorlinks=true,
  linkcolor=Blue,
  urlcolor=Blue,
  pdftitle={コンパイラ演習 2026 WhileLang から WebAssembly への翻訳},
  pdfauthor={情報システム実験 I}
}
\pagestyle{fancy}
\fancyhf{}
\lhead{情報システム実験 I}
\rhead{コンパイラ演習、伊澤侑祐}
\cfoot{\thepage}
\renewcommand{\headrulewidth}{0.4pt}
\setlength{\parindent}{1em}
\setlength{\parskip}{0.25em}
\setlist[itemize]{leftmargin=2em,itemsep=0.2em,topsep=0.3em}
\setlist[enumerate]{leftmargin=2.2em,itemsep=0.25em,topsep=0.3em}

\titleformat{\section}{\Large\bfseries\color{Ink}}{\thesection}{0.7em}{}
\titleformat{\subsection}{\large\bfseries\color{Blue}}{\thesubsection}{0.7em}{}

\lstdefinestyle{course}{
  basicstyle=\ttfamily\small,
  keywordstyle=\bfseries\color{Blue},
  commentstyle=\color{CodeGray},
  stringstyle=\color{Orange},
  showstringspaces=false,
  columns=fullflexible,
  keepspaces=true,
  frame=single,
  rulecolor=\color{LightBlue},
  backgroundcolor=\color{SoftGray},
  xleftmargin=0.6em,
  xrightmargin=0.6em,
  aboveskip=0.45em,
  belowskip=0.45em
}
\lstset{style=course}

\newcommand{\term}[1]{\textbf{\color{Blue}#1}}
\newcommand{\task}[1]{\textbf{\color{Orange}#1}}

\begin{document}

\begin{center}
  {\LARGE\bfseries\color{Ink} コンパイラ演習 2026}\par
  \vspace{0.25em}
  {\large\color{Blue} WhileLang から WebAssembly への翻訳}\par
  \vspace{0.45em}
  {\small 情報システム実験 I　集中講義用レジュメ}
\end{center}
\vspace{0.4em}
\noindent\color{Blue}\rule{\textwidth}{0.8pt}\color{Ink}

この演習では、小さな命令型言語 WhileLang を WebAssembly テキスト形式 \texttt{.wat} へ翻訳するコンパイラを実装します。
字句解析器と構文解析器は配布済みであり、構文木から仮想スタック命令への翻訳と、その命令列から WAT への出力を完成させる。目標は、コンパイラの各段階を見通しながら、算術演算、比較演算、while 文の意味を保つ翻訳を実装することである。

\section{コンパイラの見取り図}

コンパイラは、ある言語で書いたプログラムを別の言語へ翻訳する。翻訳前後で表現は変わるが、プログラムを実行した結果は同じでなければならない。今回の処理は次の五段階で進む。

\begin{center}
\begin{tikzpicture}[node distance=1cm,>=Latex]
  \tikzset{box/.style={draw=Blue,rounded corners=2pt,fill=LightBlue,
    minimum width=2.75cm,minimum height=0.82cm,align=center,font=\small}}
  \node[box] (src) {WhileLang ソース};
  \node[box,right=of src] (tok) {トークン列};
  \node[box,right=of tok] (ast) {構文木 AST};
  \node[box,below=0.65cm of ast] (stack) {仮想スタック命令};
  \node[box,left=of stack] (wat) {WebAssembly WAT};
  \draw[->,thick,Blue] (src) -- node[above,font=\scriptsize]{字句解析} (tok);
  \draw[->,thick,Blue] (tok) -- node[above,font=\scriptsize]{構文解析} (ast);
  \draw[->,thick,Blue] (ast) -- node[right,font=\scriptsize]{課題} (stack);
  \draw[->,thick,Blue] (stack) -- node[above,font=\scriptsize]{課題} (wat);
\end{tikzpicture}
\end{center}

変更するファイルは次の二つだけである。
\begin{itemize}
  \item Python: \path{virtual_stack.py} と \path{emit_wasm.py}
  \item C: \path{virtual_stack.c} と \path{emit_wasm.c}
\end{itemize}

\textbf{最初に Python または C の一方を選んだら、途中で実装を混ぜない。}

\section{仮想スタック機械}

仮想スタック命令は、値を積む、値を取り出す、二つの値を演算する、といった小さな操作からなる。
二項演算では、左辺を翻訳し、右辺を翻訳し、最後に演算命令を置く。この並びは後置記法に相当する。

\begin{minipage}[t]{0.47\textwidth}
\textbf{WhileLang}
\begin{lstlisting}
a := a * 2;
\end{lstlisting}
\end{minipage}\hfill
\begin{minipage}[t]{0.49\textwidth}
\textbf{仮想スタック命令}
\begin{lstlisting}
rvalue a
push   2
*
lpush  a
\end{lstlisting}
\end{minipage}

\noindent\textbf{重要な区別}
\begin{itemize}
  \item \texttt{rvalue a} は変数 \texttt{a} の値を積む。
  \item \texttt{lpush a} はスタック先頭の値を変数 \texttt{a} に格納する。
  \item \texttt{-} では、下にある値を左辺、上にある値を右辺として計算する。
\end{itemize}

Wasm もスタック機械なので、仮想命令の多くは WAT の命令と一対一で対応する。

\begin{center}
\begin{tabular}{lll}
\toprule
役割 & 仮想命令 & WAT \\
\midrule
定数を積む & \texttt{push n} & \texttt{i32.const n} \\
変数を読む & \texttt{rvalue x} & \texttt{global.get \$x} \\
変数へ書く & \texttt{lpush x} & \texttt{global.set \$x} \\
加算 & \texttt{+} & \texttt{i32.add} \\
符号付き比較 & \texttt{<} & \texttt{i32.lt\_s} \\
出力 & \texttt{print} & \texttt{call \$print} \\
\bottomrule
\end{tabular}
\end{center}

\section{課題と実装のヒント}

\subsection{課題1 算術演算}

\task{実装内容}は、\texttt{Sub / Mul / Div} の構文木を仮想命令へ翻訳することと、
\texttt{Minus / Times} に対して \texttt{i32.sub / i32.mul} を出力することである。割り算の WAT 出力 \texttt{i32.div\_s} は配布済みである。

\begin{lstlisting}[language=Python]
case Add(lhs, rhs):
    return compile_arith(lhs) + compile_arith(rhs) + [Plus()]
\end{lstlisting}

\texttt{Add} の場合を写し、末尾の命令を変えればよい。
C では \texttt{ARITH\_ADD} と補助関数 \texttt{bin\_arith} を手本にする。
引き算 \texttt{10 - 3} では 10、3 の順に積む。
\texttt{print 3 - 10;} が \(-7\) になれば順序は正しい。

\subsection{課題2 比較演算}

\task{実装内容}は、\texttt{GT / GE / LE / EQ} を仮想命令へ翻訳し、対応する WAT を出力することである。\texttt{LT} が完成例になっている。

\begin{center}
\begin{tabular}{cc@{\qquad}cc}
\toprule
比較 & WAT & 比較 & WAT \\
\midrule
\texttt{>} & \texttt{i32.gt\_s} & \texttt{>=} & \texttt{i32.ge\_s} \\
\texttt{<=} & \texttt{i32.le\_s} & \texttt{==} & \texttt{i32.eq} \\
\bottomrule
\end{tabular}
\end{center}

比較結果は真なら \texttt{1}、偽なら \texttt{0} になる。\texttt{print i > 1;} のように、比較式を直接出力して確認できる。

\subsection{課題3 while}

while 文には、条件を調べる場所 \texttt{test} と、ループを抜けた後の場所 \texttt{out} が必要である。

\begin{minipage}[t]{0.36\textwidth}
\begin{lstlisting}
while i < 10 do
  i := i + 1;
\end{lstlisting}
\end{minipage}\hfill
\begin{minipage}[t]{0.60\textwidth}
\begin{lstlisting}
label   L.0       # test
rvalue i
push    10
<
gofalse L.1      # out
rvalue i
push    1
+
lpush   i
goto    L.0
label   L.1
\end{lstlisting}
\end{minipage}

命令を組み立てる順序は、\texttt{LabelTest}、条件、\texttt{GoFalse}、本体、\texttt{GoTo}、\texttt{LabelOut} である。\texttt{GoFalse} は条件が 0 のとき \texttt{out} へ進み、\texttt{GoTo} は無条件に \texttt{test} へ戻る。ラベル名は各 while ごとに \texttt{gen\_label()} で新しく作れば、入れ子にも対応できる。

WAT では \texttt{block} が脱出先を、内側の \texttt{loop} が反復先を表す。ここを出力する処理は配布済みである。

\begin{lstlisting}
(block $L.1
  (loop $L.0
    ... condition ...
    i32.eqz
    br_if $L.1
    ... body ...
    br $L.0
  )
)
\end{lstlisting}

\noindent\textbf{よくある失敗}
\begin{itemize}
  \item \texttt{GoFalse} と \texttt{GoTo} を逆にして、一度も回らない、または無限ループになる。
  \item \texttt{LabelTest(test, out)} と \texttt{LabelOut(test, out)} に別の組を渡す。
  \item while の本体が複数文なのに、入力側で \texttt{begin ... end} を使わない。
\end{itemize}

\section{デバッグと確認手順}

不具合は、翻訳パイプラインの手前から順に調べる。AST が期待どおりならパーサは正しく、翻訳側を調べればよい。WAT を調べる前に \texttt{whilei} で仮想命令が正しく動くことを確認する。

\begin{enumerate}
  \item 構文木を見る。
\begin{lstlisting}
python3 whilec.py --ast ../test/arith.while
\end{lstlisting}
  \item 仮想命令列を見る。
\begin{lstlisting}
python3 whilec.py --stack ../test/arith.while
\end{lstlisting}
  \item 一命令ずつスタックと変数を追う。
\begin{lstlisting}
python3 whilei.py --trace ../test/arith.while
\end{lstlisting}
  \item 全テストを実行する。
\begin{lstlisting}
python3 test_whilelang.py
\end{lstlisting}
\end{enumerate}

C 版では \texttt{./whilec}、\texttt{./whilei} に読み替え、全テストは \texttt{make test} で実行する。Windows では Python を \texttt{py -3} で起動できる。

\begin{center}
\begin{tabularx}{0.88\textwidth}{>{\ttfamily}l>{\ttfamily}lX}
\toprule
入力 & 期待する出力 & 確認する内容 \\
\midrule
assign.while & 3 & 配布状態と加算 \\
arith.while & 7 8 4 23 & 減算、乗算、除算、優先順位 \\
cmp.while & 1 1 0 1 1 0 & 比較結果 \\
simple\_loop.while & 10 & 基本の while \\
fact.while & 120 & while と乗算 \\
loop.while & 入れ子の出力 & ラベルの組 \\
\bottomrule
\end{tabularx}
\end{center}

\section{着手順と発展課題}

最初に配布状態で \texttt{assign.while} が \texttt{3} を出すことを確かめる。その後、課題1、課題2、課題3の順に進み、短いテストを通してから次へ移る。詰まったときはリポジトリ直下の \texttt{HINTS.md} を上から読む。

時間が余ったら、\texttt{if P then S else S} を実装する。命令の順序は、条件、\texttt{IfStart}、then 節、\texttt{ElseOp}、else 節、\texttt{IfEnd} である。Python の命令名は \texttt{IfStart / ElseOp / IfEnd}、C では \texttt{I\_IF\_START / I\_ELSE / I\_IF\_END} である。

\vfill
\noindent\small\textbf{参照先}　\texttt{README.md}、\texttt{python/README.md}、\texttt{c/README.md}、\texttt{HINTS.md}

\end{document}
